\documentclass{article}
\usepackage[utf8]{inputenc}
\usepackage[T1]{fontenc}
\usepackage{geometry}
\usepackage{graphicx}
\usepackage{float}
\usepackage{hyperref}
\usepackage[english,indonesian]{babel}

\title{Optimasi Pemrosesan Template LaTeX Menggunakan Pendekatan Berbasis AI}
\author{Academic Support Agent}
\date{\today}

\begin{document}
\maketitle

\begin{abstract}

\end{abstract}

\chapter{PENDAHULUAN}


\chapter{TINJAUAN PUSTAKA}


% Original content (commented out):
% %Aquest document és una traducció i adaptació al català
%del "sample paper template" proporcionat per Overleaf

\documentclass[a4paper]{article}

%% idioma i codificació 
\usepackage[catalan]{babel}
\usepackage[utf8x]{inputenc}
\usepackage[T1]{fontenc}

%% mida de pàgina i marges
\usepackage[a4paper,top=3cm,bottom=2cm,left=3cm,right=3cm,marginparwidth=1.75cm]{geometry}

%% alguns paquets útils
\usepackage{amsmath} % per fórmules
\usepackage{graphicx} % permet inserir imatges en el nostre document
\usepackage[colorlinks=true, allcolors=blue]{hyperref} % hiperenllaços
\usepackage{booktabs} % per fer taules més elegants


\title{Breu introducció a \LaTeX}
\author{Ignasi Furió}

\newcommand{\comanda}[1][]{\texttt{\textbackslash{}#1}}

\begin{document}
\maketitle

\begin{abstract}
En aquesta sessió intentarem presentar les comandes bàsiques de LaTeX que més útils vos poden ser per poder presentar treballs en aquest sistema.
\end{abstract}

\section{Introducció}

Per començar podeu observar que podem escriure el text normalment sense preocupar-nos de gaires coses més. El sistema divideix les línies i paràgrafs automàticament, fent que s'ajustin als marges indicats.

Si voleu començar un nou paràgraf, deixau una línia en blanc entre un i l'altre.

Com podeu observar       els         espais en blanc
o fins i tot les noves línies 
no signifiquen res a efectes del resultat final.

Dit això, començarem amb una sèrie de comandes bàsiques.

\section{Alguns exemples per començar}

\subsection{Com incloure figures}\label{sec:imatges}

Si treballam amb Overleaf, pujarem la imatge des del nostre ordinador a la carpeta del projecte mitjançant l'opció \emph{upload} del menú de projecte. Desprès usarem la comanda \texttt{\textbackslash{}includegraphics} per incloure-la en el nostre document. 

Molt important les figures són objectes flotants, per tant, el sistema les posarà on convengui. No patiu, com que hi posarem un peu i una referència, sempre podrem esmentar-les dins el text sense cap problema. Per exemple, aquí mateix podem fer referència a la imatge de la granota (figura~\ref{fig:frog}) que apareix en algun lloc d'aquest document. El codi per crear la figura és el que apareix a continuació.

\begin{figure}
\centering
\includegraphics[width=0.3\textwidth]{frog}
\caption{\label{fig:frog}Això és una imatge d'una granota que tenim en el nostre projecte.}
\end{figure}

\subsection{Com fer taules}

En \LaTeX{} heu de distingir entre una \emph{table} i un \emph{tabular}. Les primeres són objectes flotants, com les figures, el sistema decideix quina és la millor posició per co"locar-la, solen tenir un peu i una etiqueta per poder fer-ne referència. Els \emph{tabular} són els entorns que ens permetin distribuir dades en files i columnes i que solen esser els continguts de les \emph{tables}.  

Dins el fitxer \texttt{.tex} podreu llegir a continuació d'aquesta frase com fer una taula senzilla com la que es mostra a \ref{tab:widgets}.

\begin{table}
\centering
\begin{tabular}{l|r}
Item & Quantitat \\\hline
Widgets & 42 \\
Gadgets & 13
\end{tabular}
\caption{\label{tab:widgets}Un exemple de taula.}
\end{table}

Encara que el format anterior el trobareu sovint en molts documents,  les normes tipogràfiques recomanen no utilitzar mai línies verticals dins les taules. Una forma més elegant de realitzar la mateixa taula pot ser la següent (amb l'ajuda del paquet \emph{booktabs}). El resultat el trobareu a la taula~\ref{tab:widgets2}.

\begin{table}
\centering
\begin{tabular}{lr}
\toprule
\multicolumn{1}{c}{\textbf{Item}} & \multicolumn{1}{c}{\textbf{Quantitat}}\\
\midrule
Widgets & 42 \\
Gadgets & 13 \\
\bottomrule
\end{tabular}
\caption{\label{tab:widgets2}Un exemple de taula sense línies verticals.}
\end{table}

\subsection{Com escriure matemàtiques}

\LaTeX{} és fantàstic per escriure fórmules matemàtiques. Un exemple senzill: Sigui $X_1, X_2, \ldots, X_n$ una seqüència de variables aleatòries independents i idènticament distribuïdes amb $\text{E}[X_i] = \mu$ i $\text{Var}[X_i] = \sigma^2 < \infty$, i sigui
\[S_n = \frac{X_1 + X_2 + \cdots + X_n}{n}
      = \frac{1}{n}\sum_{i}^{n} X_i\]
la seva mitjana. Aleshores, quan $n$ s'apropa a infinit, les variables aleatòries $\sqrt{n}(S_n - \mu)$ convergeixen a una distribució normal $\mathcal{N}(0, \sigma^2)$.

\subsection{Com crear seccions i subseccions}

L'ús de seccions i subseccions (i altres divisions) permet organitzar el nostre document. Sols cal incloure el nom corresponent dins la comanda. \LaTeX{}  s'encarregarà de crear la divisió dins el text i afegir el format i numeració que pertoqui. També podem afegir etiquetes per tal de fer-ne referència quan convengui. Per exemple, hem explicat com incloure imatges a la secció~\ref{sec:imatges}.

Les divisions disponibles són: \comanda{part}, \comanda{chapter}\footnote{No disponible dins \texttt{article}}, \comanda{section}, \comanda{subsection}, \comanda{subsubsection}, \comanda{paragraph}, \comanda{subparagraph}.

\subsection{Com fer llistes}

Podem fer llistes amb numeració automàtica 

\begin{enumerate}
\item Aquest és un ítem,
\item i aquest un altre.
\end{enumerate}

o sense enumerar

\begin{itemize}
\item Aquest és un ítem,
\item i aquest un altre.
\end{itemize}

\subsection{Com afegir cites i llistes de referències}

Podem pujar a la nostra carpeta un fitxer \verb|.bib| que conté les nostres referències creat a ma o amb eines com JabRef o importar el fitxer des de serveis de gestió bibliogràfica com \href{https://www.overleaf.com/blog/184}{Mendeley}, CiteULike or Zotero. Podem incloure les cites de les referències que en formen part, com per exemple aquesta: \cite{greenwade93}. Recorda que has d'especificar l'estil bibliogràfic que vols usar i el nom del fitxer \verb|.bib|.

Podeu mirar \href{https://www.overleaf.com/help/97-how-to-include-a-bibliography-using-bibtex}{aquest video tutorial} per aprendre una mica més de BibTeX.

\bibliographystyle{alpha}
\bibliography{sample}

\end{document}

\end{document}